% !TEX root = ../../ctfp-print.tex

\lettrine[lhang=0.17]{范}{畴论中的大多数构造}都是其他更具体数学分支中结果的推广。诸如积、余积、幺半群、指数对象等概念早在范畴论诞生之前就已存在。它们可能在不同数学分支中有不同的名称。集合论中的笛卡尔积、序理论中的交运算、逻辑学中的合取——这些都是范畴积抽象概念的具体实例。

在这方面，Yoneda引理（Yoneda lemma）显得尤为突出，它是一个关于范畴的普遍性陈述，在其他数学分支中几乎没有先例。有人认为其最接近的类比是群论中的凯莱定理（每个群都同构于某个集合的置换群）。

Yoneda引理的背景是一个任意范畴$\cat{C}$以及一个从$\cat{C}$到$\Set$的函子$F$。我们在上一节见过，某些$\Set$-值函子是可表的，即同构于某个Hom函子。Yoneda引理告诉我们，所有$\Set$-值函子都可以通过自然变换从Hom函子得到，并且明确列举了所有此类变换。

当我谈论自然变换时，曾提到自然性条件可能非常强约束。当你在一个对象上定义自然变换的一个分量时，自然性可能足以将这个分量"传输"到通过态射连接的另一个对象。源范畴和目标范畴中对象间的箭头越多，自然变换分量传输时的约束就越强。$\Set$恰好是一个具有丰富态射的范畴。

Yoneda引理告诉我们，Hom函子与任意其他函子$F$之间的自然变换，完全由指定其单个分量在一个点上的值所确定！其余部分仅由自然性条件推导得出。

让我们回顾Yoneda引理涉及的两个函子之间的自然性条件。第一个函子是Hom函子，它将$\cat{C}$中任意对象$x$映射到态射集$\cat{C}(a, x)$（其中$a$是$\cat{C}$中的固定对象）。我们也见过它将任意态射$f:x \to y$映射到$\cat{C}(a, f)$。

第二个函子是任意的$\Set$-值函子$F$。

设这两个函子之间的自然变换为$\alpha$。由于我们在$\Set$中操作，自然变换的分量（如$\alpha_x$或$\alpha_y$）只是集合间的常规函数：
\begin{gather*}
  \alpha_x \Colon \cat{C}(a, x) \to F x \\
  \alpha_y \Colon \cat{C}(a, y) \to F y
\end{gather*}

\begin{figure}[H]
  \centering
  \includegraphics[width=0.4\textwidth]{images/yoneda1.png}
\end{figure}

\noindent
因为这些只是函数，我们可以考察它们在特定点处的值。但集合$\cat{C}(a, x)$中的"点"是什么？关键观察在于：集合$\cat{C}(a, x)$中的每个点也是一个从$a$到$x$的态射$h$。

因此，$\alpha$的自然性方阵：
$$\alpha_y \circ \cat{C}(a, f) = F f \circ \alpha_x$$
在作用于$h$时变为逐点形式：
$$\alpha_y (\cat{C}(a, f) h) = (F f) (\alpha_x h)$$
回想上一节中Hom函子$\cat{C}(a,-)$作用于态射$f$的定义是前组合：
$$\cat{C}(a, f) h = f \circ h$$
代入得：
$$\alpha_y (f \circ h) = (F f) (\alpha_x h)$$
当特化到$x = a$情形时，这种条件的强大程度变得明显：

\begin{figure}[H]
  \centering
  \includegraphics[width=0.4\textwidth]{images/yoneda2.png}
\end{figure}

\noindent
此时$h$成为从$a$到$a$的态射。我们知道至少存在一个这样的态射$h = \id_a$。将其代入：
$$\alpha_y f = (F f) (\alpha_a \id_a)$$
注意发生了什么：左侧是$\alpha_y$作用于$\cat{C}(a, y)$中任意元素$f$的结果，而这完全由$\alpha_a$在$\id_a$处的单个值确定。我们可以任选这样的值，它都会生成一个自然变换。由于$\alpha_a$的值在集合$F a$中，$F a$中的任何点都将定义某个$\alpha$。

反之，给定从$\cat{C}(a, -)$到$F$的任意自然变换$\alpha$，你可以在$\id_a$处求值得到$F a$中的一个点。

我们刚刚证明了Yoneda引理：

\begin{quote}
  从$\cat{C}(a, -)$到$F$的自然变换与$F a$的元素之间存在一一对应关系。
\end{quote}
换句话说，
$$\mathit{Nat}(\cat{C}(a, -), F) \cong F a$$
或者，如果我们用$[\cat{C}, \Set]$表示$\cat{C}$到$\Set$的函子范畴，则自然变换集正是该范畴中的Hom集，可以写作：
$$[\cat{C}, \Set](\cat{C}(a, -), F) \cong F a$$
稍后我将解释这种对应实际上是如何构成自然同构的。

现在尝试获得对该结果的直观理解。最惊人的事情是整个自然变换从一个单一的成核点结晶而出：我们在$\id_a$处分配的值。它遵循自然性条件从该点扩散开来，淹没$\cat{C}$在$\Set$中的像。因此让我们首先考虑$\cat{C}(a, -)$下$\cat{C}$的像是什么。

从$a$自身的像开始：在Hom函子$\cat{C}(a, -)$下，$a$被映射到集合$\cat{C}(a, a)$。而在函子$F$下，它被映射到集合$F a$。自然变换$\alpha_a$的分量是从$\cat{C}(a, a)$到$F a$的某个函数。让我们聚焦于集合$\cat{C}(a, a)$中对应态射$\id_a$的那个点。为了强调这只是集合中的一个点，我们称它为$p$。分量$\alpha_a$应该将$p$映射到$F a$中的某个点$q$。我将向你展示任何$q$的选择都会导致唯一的自然变换。

\begin{figure}[H]
  \centering
  \includegraphics[width=0.3\textwidth]{images/yoneda3.png}
\end{figure}

\noindent
第一个断言是：选择单个点$q$唯一确定了函数$\alpha_a$的其余部分。确实，让我们选取$\cat{C}(a, a)$中任意其他点$p'$，对应某个从$a$到$a$的态射$g$。这就是Yoneda引理魔力发生的地方：$g$可以被视为集合$\cat{C}(a, a)$中的点$p'$。同时，它选择了两个\emph{函数}在集合之间。事实上，在Hom函子下，态射$g$被映射到函数$\cat{C}(a, g)$；而在$F$下它被映射到$F g$。

\begin{figure}[H]
  \centering
  \includegraphics[width=0.4\textwidth]{images/yoneda4.png}
\end{figure}

\noindent
现在考虑$\cat{C}(a, g)$对我们原始的$p$（记得对应$\id_a$）的作用。其定义为前组合$g \circ \id_a$，等于$g$，对应我们的点$p'$。因此态射$g$被映射到这样一个函数：当作用于$p$时产生$p'$（即$g$）。我们已经完成了一个循环！

现在考虑$F g$作用于$q$。这是$F a$中的某个点$q'$。为了完成自然性方阵，$p'$必须在$\alpha_a$下被映射到$q'$。我们选取了任意的$p'$（任意的$g$），并推导出它在$\alpha_a$下的映射。因此函数$\alpha_a$被完全确定。

第二个断言是：对于$\cat{C}$中任何与$a$相连的对象$x$，$\alpha_x$都是唯一确定的。推理过程类似，只不过现在我们有两个更多集合$\cat{C}(a, x)$和$F x$，态射$g:a \to x$在Hom函子下被映射为：
$$\cat{C}(a, g) \Colon \cat{C}(a, a) \to \cat{C}(a, x)$$
而在$F$下被映射为：
$$F g \Colon F a \to F x$$
同样，$\cat{C}(a, g)$作用于我们的$p$由前组合给出：$g \circ \id_a$，对应$\cat{C}(a, x)$中的点$p'$。自然性确定了$\alpha_x$作用于$p'$的值为：
$$q' = (F g) q$$
由于$p'$是任意的，整个函数$\alpha_x$因此被确定。

\begin{figure}[H]
  \centering
  \includegraphics[width=0.4\textwidth]{images/yoneda5.png}
\end{figure}

\noindent
如果$\cat{C}$中存在与$a$没有连接的对象怎么办？它们都在$\cat{C}(a, -)$下被映射到单个集合——空集。回忆空集是集合范畴的始对象，这意味着从这个集合到任何其他集合都有唯一函数。我们将这个函数称为\code{absurd}。因此在这种情况下，自然变换的分量别无选择：只能是\code{absurd}。

理解Yoneda引理的一种方式是认识到$\Set$-值函子之间的自然变换只是函数族，而函数通常是信息有损的。函数可能压缩信息，也可能仅覆盖其陪域的一部分。唯一无损的函数是可逆函数——同构。因此最好的保结构$\Set$-值函子是可表的函子。它们要么是Hom函子，要么是与Hom函子自然同构的函子。任何其他函子$F$都是通过有损变换从Hom函子得到的。这样的变换不仅可能丢失信息，还可能仅覆盖函子$F$在$\Set$中像的一小部分。

\section{Haskell中的Yoneda}

我们已经在Haskell中以读者函子的形式遇到过Hom函子：

\src{snippet01}
读者函子通过前组合映射态射（在这里是函数）：

\src{snippet02}
Yoneda引理告诉我们读者函子可以自然地映射到任何其他函子。

自然变换是一个多态函数。因此给定函子\code{F}，我们有一个从读者函子到它的映射：

\src{snippet03}
像往常一样，\code{forall}是可选的，但我喜欢显式写出它以强调自然变换的参数多态性。

Yoneda引理告诉我们这些自然变换与\code{F a}的元素一一对应：

\begin{snipv}
forall x . (a -> x) -> F x \ensuremath{\cong} F a
\end{snipv}
这个恒等式的右侧是我们通常会视为数据结构的东西。记住函子作为广义容器的解释？\code{F a}是\code{a}的容器。但左侧是一个接受函数作为参数的多态函数。Yoneda引理告诉我们这两种表示是等价的——它们包含相同的信息。

另一种说法是：给我一个如下类型的多态函数：

\src{snippet04}
我就能生成一个\code{a}的容器。诀窍就是我们在Yoneda引理证明中使用的技巧：用\code{id}调用这个函数得到\code{F a}的元素：

\src{snippet05}
反过来也成立：给定类型为\code{F a}的值：

\src{snippet06}
可以定义一个如下类型的多态函数：

\src{snippet07}
具有正确的类型。你可以轻松在两种表示之间来回转换。

拥有多种表示的优势在于其中一种可能比另一种更容易组合，或者在某些应用中更高效。

这一原理最简单的例子是编译器构造中常用的代码转换技术：延续传递风格（Continuation Passing Style，简称\acronym{CPS}）。它是Yoneda引理应用于恒等函子的最简单应用。将\code{F}替换为恒等函子得到：

\begin{snipv}
forall r . (a -> r) -> r \ensuremath{\cong} a
\end{snipv}
这个公式的解释是：任何类型\code{a}都可以被一个接受\code{a}的"处理器"（continuation handler）的函数替代。处理器是一个接受\code{a}并执行剩余计算的函数——即延续（continuation）。（类型\code{r}通常封装某种状态码。）

这种编程风格在UI开发、异步系统和并发编程中非常常见。\acronym{CPS}的缺点是涉及控制反转。代码被分割为生产者和消费者（处理器），难以组合。任何做过大量非平凡Web编程的人都熟悉交互式有状态处理器带来的意大利面条式代码噩梦。正如我们稍后将看到的，明智地使用函子和单子可以恢复\acronym{CPS}的一些组合性质。

\section{共Yoneda}

与往常一样，通过反转箭头的方向我们可以得到一个额外构造。Yoneda引理可应用于相反范畴$\cat{C}^\mathit{op}$，从而在逆变函子之间建立映射。

等价地，我们也可以通过固定同态函子的目标对象而非源对象来推导共Yoneda引理。我们得到从范畴$\cat{C}$到$\Set$的逆变同态函子：
$\cat{C}(-, a)$。Yoneda引理的逆变版本建立了从该函子到任意其他逆变函子 $F$ 的自然变换与集合 $F a$ 元素之间的一一对应关系：
$$\mathit{Nat}(\cat{C}(-, a), F) \cong F a$$
以下是共Yoneda引理的Haskell版本：

\begin{snipv}
forall x . (x -> a) -> F x \ensuremath{\cong} F a
\end{snipv}
请注意，在某些文献中，这个逆变版本被称为Yoneda引理。

\section{挑战}

\begin{enumerate}
  \tightlist
  \item
        证明构成Haskell中Yoneda同构的两个函数 \code{phi} 和 \code{psi} 互为逆元。

        \begin{snip}{haskell}
phi :: (forall x . (a -> x) -> F x) -> F a
phi alpha = alpha id

psi :: F a -> (forall x . (a -> x) -> F x)
psi fa h = fmap h fa
\end{snip}
  \item
        离散范畴(discrete category)是指仅有对象而没有恒等态射(identity morphism)之外其他态射(morphism)的范畴。这种情况下Yoneda引理对这类范畴上的函子(functor)如何成立？
  \item
        单位元列表 \code{{[}(){]}} 除了其长度外不含任何其他信息。因此作为数据类型(data type)，它可以被视为整数的编码方式。空列表编码零，单元素列表 \code{{[}(){]}}（值而非类型）编码一，依此类推。请使用列表函子(list functor)的Yoneda引理构造该数据类型的另一种表示形式。
\end{enumerate}

\section{参考文献}

\begin{enumerate}
  \tightlist
  \item
        \urlref{https://www.youtube.com/watch?v=TLMxHB19khE}{Catsters 视频}
\end{enumerate}